\section{Modelling of Digital Control Systems}\label{sec:modelling}
\lecdate{14 and 21 October 2019}
% Block diagrams of Chapter 3 (redrawn from the board sketches).
% Needs: dtcnotes.sty (styles blk, sumj, sig, lab, dot and the sampler pic).

% ---- p.26: physical digital control system ---------------------------------------
\newcommand{\figDigitalLoop}{%
\begin{tikzpicture}
  \node[blk,minimum width=33mm,minimum height=11mm] (cpu) {Computer or\\ microprocessor};
  \node[blk,right=10mm of cpu] (dac) {DAC};
  \node[blk,right=10mm of dac] (as) {Analog\\ system};
  \node[blk,below=13mm of dac] (adc) {ADC};
  \draw[sig] ([xshift=-12mm]cpu.west) -- (cpu.west) node[pos=0,above right,lab]{input};
  \draw[sig] (cpu) -- (dac);
  \draw[sig] (dac) -- (as);
  \coordinate (br) at ($(as.east)+(10mm,0)$);
  \draw[sig] (as.east) -- (br) -- ++(8mm,0) node[right,lab,align=left]{Analog\\ output};
  \fill (br) circle (1.4pt);
  \draw[sig] (br) |- (adc.east);
  \draw[sig] (adc.west) -| (cpu.south);
\end{tikzpicture}}

% ---- p.26: equivalent discrete-time loop -----------------------------------------
\newcommand{\figEquivLoop}[2]{% #1 = reference label, #2 = output label
\begin{tikzpicture}
  \node[sumj] (s) {};
  \node[blk,right=13mm of s] (c) {$C(z)$};
  \node[blk,right=15mm of c] (g) {$\GZAS(z)$};
  \draw[sig] ([xshift=-14mm]s.west) -- (s) node[pos=0,above right,lab]{#1};
  \node[font=\scriptsize] at ($(s.west)+(-2.5mm,2.6mm)$) {$+$};
  \node[font=\scriptsize] at ($(s.south)+(-2.6mm,-2.4mm)$) {$-$};
  \draw[sig] (s) -- (c) node[midway,above,lab]{$E(z)$};
  \draw[sig] (c) -- (g);
  \coordinate (br) at ($(g.east)+(12mm,0)$);
  \draw[sig] (g.east) -- (br) -- ++(6mm,0) node[right,lab]{#2};
  \fill (br) circle (1.4pt);
  \draw[sig] (br) -- ++(0,-14mm) -| (s.south);
\end{tikzpicture}}

% ---- p.26: what G_ZAS(z) contains -----------------------------------------------
\newcommand{\figZASdef}{%
\begin{tikzpicture}
  \node[blk] (zoh) {ZOH};
  \node[blk,right=10mm of zoh] (as) {Analog\\ subsys.};
  \pic (sm) at ([xshift=9mm]as.east) {sampler};
  \draw[sig] ([xshift=-10mm]zoh.west) -- (zoh);
  \draw[sig] (zoh) -- (as);
  \draw[sig] (as.east) -- (sm-in);
  \draw[sig] (sm-out) -- ++(10mm,0);
  \node[lab,above=1mm of zoh] {DAC};
  \node[lab,below=4mm of sm-in,xshift=6mm] {ideal sampler (ADC)};
\end{tikzpicture}}

% ---- p.34: starred signals --------------------------------------------------------
\newcommand{\figStarred}{%
\begin{tikzpicture}
  \pic (s1) at (0,0) {sampler};
  \node[blk,right=18mm of s1-out] (g) {$G(s)$};
  \pic (s2) at ([xshift=18mm]g.east) {sampler};
  \draw[sig] ([xshift=-14mm]s1-in) -- (s1-in) node[pos=0.45,above,lab]{$x(t)$} node[pos=0.45,below,lab]{$X(s)$};
  \draw[sig] (s1-out) -- (g) node[midway,above,lab]{$x^{*}(t)$} node[midway,below,lab]{$X^{*}(s)$};
  \draw[sig] (g) -- (s2-in) node[midway,above,lab]{$y(t)$} node[midway,below,lab]{$Y(s)$};
  \draw[sig] (s2-out) -- ++(20mm,0) node[midway,above,lab]{$y^{*}(t)$} node[midway,below,lab]{$Y^{*}(s)$};
\end{tikzpicture}}

% ---- p.36/37: cascade NOT separated by a sampler ---------------------------------
\newcommand{\figCascadeNS}{%
\begin{tikzpicture}
  \pic (s1) at (1.4,0) {sampler};
  \node[blk,minimum width=11mm] (h1) at (4.3,0) {$H_1(s)$};
  \node[blk,minimum width=11mm] (h2) at (6.6,0) {$H_2(s)$};
  \pic (s3) at (8.6,0) {sampler};
  \draw[sig] (0,0) -- (s1-in) node[pos=0.5,above,lab]{$U(s)$};
  \draw[sig] (s1-out) -- (h1) node[midway,above,lab]{$U^{*}(s)$};
  \draw[sig] (h1) -- (h2) node[midway,above,lab]{$X(s)$};
  \draw[sig] (h2) -- (s3-in) node[midway,above,lab]{$Y(s)$};
  \draw[sig] (s3-out) -- ++(17mm,0) node[midway,above,lab]{$Y^{*}(s)$};
\end{tikzpicture}}

% ---- p.36/37: cascade separated by a sampler -------------------------------------
\newcommand{\figCascadeS}{%
\begin{tikzpicture}
  \pic (s1) at (1.4,0) {sampler};
  \node[blk,minimum width=11mm] (h1) at (4.3,0) {$H_1(s)$};
  \pic (s2) at (6.1,0) {sampler};
  \node[blk,minimum width=11mm] (h2) at (8.8,0) {$H_2(s)$};
  \pic (s3) at (10.8,0) {sampler};
  \draw[sig] (0,0) -- (s1-in) node[pos=0.5,above,lab]{$U(s)$};
  \draw[sig] (s1-out) -- (h1) node[midway,above,lab]{$U^{*}(s)$};
  \draw[sig] (h1) -- (s2-in) node[midway,above,lab]{$X(s)$};
  \draw[sig] (s2-out) -- (h2) node[midway,above,lab]{$X^{*}(s)$};
  \draw[sig] (h2) -- (s3-in) node[midway,above,lab]{$Y(s)$};
  \draw[sig] (s3-out) -- ++(17mm,0) node[midway,above,lab]{$Y^{*}(s)$};
\end{tikzpicture}}

% ---- p.38: closed loop with G_ZAS expanded ---------------------------------------
\newcommand{\figClosedLoopExpanded}{%
\begin{tikzpicture}
  \node[sumj] (s) {};
  \node[blk,right=12mm of s] (c) {$C(z)$};
  \node[blk,right=24mm of c] (zoh) {$G_{\mathrm{ZOH}}(s)$};
  \node[blk,right=9mm of zoh] (g) {$G(s)$};
  \pic (sm) at ([xshift=8mm]g.east) {sampler};
  \node[draw=itublue,thick,rounded corners,fit=(zoh)(g)(sm-in)(sm-out),inner sep=7pt,
        label={[lab,anchor=south east]north east:$\GZAS(z)$}] (box) {};
  \draw[sig] ([xshift=-14mm]s.west) -- (s) node[pos=0,above right,lab]{$R(z)$};
  \node[font=\scriptsize] at ($(s.west)+(-2.5mm,2.6mm)$) {$+$};
  \node[font=\scriptsize] at ($(s.south)+(-2.6mm,-2.4mm)$) {$-$};
  \draw[sig] (s) -- (c) node[midway,above,lab]{$E(z)$};
  \draw[sig] (c) -- (box.west) node[midway,above,lab]{$U(z)$};
  \draw[sig] (box.west) -- (zoh);
  \draw[sig] (zoh) -- (g);
  \draw[sig] (g) -- (sm-in);
  \coordinate (br) at ($(box.east)+(10mm,0)$);
  \draw[sig] (sm-out) -- (box.east) -- (br) -- ++(6mm,0) node[right,lab]{$Y(z)$};
  \fill (br) circle (1.4pt);
  \draw[sig] (br) -- ++(0,-17mm) -| (s.south);
\end{tikzpicture}}

% ---- p.39: sampled feedback with H(s) --------------------------------------------
\newcommand{\figSampledFeedback}{%
\begin{tikzpicture}
  \node[sumj] (s) {};
  \node[blk,right=14mm of s] (d) {$D(s)$};
  \node[blk,right=14mm of d] (g) {$G(s)$};
  \draw[sig] ([xshift=-14mm]s.west) -- (s) node[pos=0,above right,lab]{$R(s)$};
  \node[font=\scriptsize] at ($(s.west)+(-2.5mm,2.6mm)$) {$+$};
  \node[font=\scriptsize] at ($(s.south)+(-2.6mm,-2.4mm)$) {$-$};
  \draw[sig] (s) -- (d) node[midway,above,lab]{$E(s)$};
  \draw[sig] (d) -- (g);
  \coordinate (br) at ($(g.east)+(14mm,0)$);
  \fill (br) circle (1.4pt);
  \draw[sig] (g) -- (br) node[midway,above,lab]{$Y(s)$};
  \pic (so) at ([xshift=12mm]br) {sampler};
  \draw[sig] (br) -- (so-in);
  \draw[sig] (so-out) -- ++(14mm,0) node[midway,above,lab]{$Y^{*}(s)$};
  % feedback path (right to left)
  \node[blk,below=15mm of g] (h) {$H(s)$};
  \pic[xscale=-1] (sf) at ($(h.west)+(-12mm,0)$) {sampler};
  \draw[sig] (br) |- (h.east) node[pos=0.75,above,lab]{$Y(s)$};
  \draw[sig] (h.west) -- (sf-in) node[midway,above,lab]{$X(s)$};
  \draw[sig] (sf-out) -| (s.south) node[pos=0.25,above,lab]{$X^{*}(s)$};
\end{tikzpicture}}

% ---- p.40: analog disturbance in a digital loop ----------------------------------
\newcommand{\figDisturbance}{%
\begin{tikzpicture}
  \node[sumj] (s) {};
  \node[blk,minimum width=12mm,right=10mm of s] (c) {$C^{*}(s)$};
  \pic (s1) at ([xshift=9mm]c.east) {sampler};
  \node[blk,right=10mm of s1-out] (zoh) {$G_{\mathrm{ZOH}}(s)$};
  \node[sumj,right=12mm of zoh] (p) {};
  \node[blk,right=12mm of p] (g) {$G(s)$};
  \pic (s2) at ([xshift=10mm]g.east) {sampler};
  \node[blk,above=10mm of p,minimum width=10mm] (gd) {$G_d(s)$};
  \draw[sig] ([xshift=-17mm]s.west) -- (s) node[pos=0,above right,lab]{$R(s)$};
  \node[font=\scriptsize] at ($(s.west)+(-2.5mm,2.6mm)$) {$+$};
  \node[font=\scriptsize] at ($(s.south)+(-2.6mm,-2.4mm)$) {$-$};
  \draw[sig] (s) -- (c) node[midway,above,lab]{$E(s)$};
  \draw[sig] (c.east) -- (s1-in) node[midway,above,lab]{$U(s)$};
  \draw[sig] (s1-out) -- (zoh) node[midway,above,lab]{$U^{*}(s)$};
  \draw[sig] (zoh) -- (p);
  \draw[sig] (p) -- (g);
  \draw[sig] ([yshift=0mm]gd.south) -- (p.north);
  \draw[sig] ([yshift=8mm]gd.north) node[above,lab]{$D(s)$} -- (gd.north);
  \draw[sig] (g) -- (s2-in) node[midway,above,lab]{$Y(s)$};
  \coordinate (br) at ($(s2-out)+(10mm,0)$);
  \draw[sig] (s2-out) -- (br) -- ++(6mm,0) node[right,lab]{$Y^{*}(s)$};
  \fill (br) circle (1.4pt);
  \draw[sig] (br) -- ++(0,-14mm) -| (s.south);
  \node[font=\scriptsize] at ($(p.west)+(-2.5mm,2.5mm)$) {$+$};
  \node[font=\scriptsize] at ($(p.north)+(-2.5mm,2.5mm)$) {$+$};
\end{tikzpicture}}


%==============================================================================
\subsection{Structure of a digital control system}\scanref{p.\,26}
In a digital control system a computer (or microprocessor) computes the control signal from the sampled
output; a DAC turns it into an analog signal for the analog system, and an ADC samples the analog output
for the computer.

\begin{figure}[ht]
\centering
\resizebox{0.86\linewidth}{!}{\figDigitalLoop}\\[4mm]
$\equiv$\\[2mm]
\resizebox{0.86\linewidth}{!}{\figEquivLoop{$R(z)$}{}}\\[4mm]
\begin{tabular}{@{}r@{\ }c@{\ }l@{}}
$\GZAS(z)$ & $\equiv$ & \raisebox{-0.9em}{\resizebox{0.5\linewidth}{!}{\figZASdef}}
\end{tabular}
\caption{A digital control system and its equivalent discrete-time representation.}\label{fig:digitalloop}
\end{figure}

\begin{itemize}
  \item $C(z)$: transfer function of the digital control system (the controller).
  \item $\GZAS(z)$: $z$ transform of the DAC (ZOH), the analog subsystem and the ADC (ideal sampler),
        taken together. (ZOH = zero-order hold, ZAS = ``zero-order hold, analog subsystem, sampler''.)
\end{itemize}

\begin{boardnote}
In the board sketch of $\GZAS(z)$ the labels ``ADC'' and ``DAC'' sit over the wrong blocks: the
ZOH models the \emph{DAC} and the ideal sampler models the \emph{ADC}. The text of the same page
($\GZAS$ contains ``DAC (ZOH), analog subsystem and ADC (ideal sampler)'') is correct, and the figure
above uses the correct labels.
\end{boardnote}

%------------------------------------------------------------------------------
\subsection{ADC model: the ideal sampler}\scanref{pp.\,26--28}
\paragraph{Assumptions.}
\begin{itemize}
  \item The ADC outputs are exactly equal in magnitude to its inputs (quantization errors are negligible).
  \item The ADC yields a digital output instantaneously.
  \item Sampling is perfectly uniform (it occurs at a fixed rate).
\end{itemize}
Under these assumptions the ADC can be modelled as an \emph{ideal sampler} with sampling period $T$.

\begin{figure}[ht]
\centering
\begin{tikzpicture}
  \begin{axis}[name=a,width=0.40\linewidth,height=3.8cm,axis lines=left,xlabel={$t$},ylabel={$u(t)$},
      xmin=0,xmax=6,ymin=0,ymax=2.5,xtick=\empty,ytick=\empty,clip=false]
    \addplot[itublue,very thick,domain=0:5.6,samples=70]{1.1+0.9*sin(deg(1.25*x+2.4))};
  \end{axis}
  \begin{axis}[at={(a.east)},anchor=west,xshift=3.4cm,width=0.44\linewidth,height=3.8cm,axis lines=left,
      xlabel={$t$},ylabel={$u^{*}(t)$},xmin=0,xmax=6,ymin=0,ymax=2.5,ytick=\empty,
      xtick={0.5,1,1.5,5},xticklabels={$T$,$2T$,$3T$,$kT$},xticklabel style={font=\scriptsize},clip=false]
    \addplot[gray,dashed,domain=0:5.6,samples=70]{1.1+0.9*sin(deg(1.25*x+2.4))};
    \addplot[ycomb,itublue,thick,mark=triangle*,mark size=2.2pt,
      samples at={0,0.5,...,5.5}]{1.1+0.9*sin(deg(1.25*x+2.4))};
  \end{axis}
  \pic at ($(a.east)+(1.3cm,0)$) {sampler};
  \draw[->,thick] ($(a.east)+(0.5cm,0)$) -- ($(a.east)+(1.3cm,0)$);
  \draw[->,thick] ($(a.east)+(2.17cm,0)$) -- ($(a.east)+(2.8cm,0)$);
\end{tikzpicture}
\caption{Ideal sampler: a continuous signal $u(t)$ becomes the impulse train $u^{*}(t)$ (impulse sampler).}
\end{figure}

The sampled signal $u^{*}(t)$, a train of impulses, is represented by
\begin{equation}
  u^{*}(t)=\sum_{k=0}^{\infty}u(kT)\,\delta(t-kT)
  =\underbrace{u(0)}_{u_0}\delta(t)+\underbrace{u(T)}_{u_1}\delta(t-T)+\underbrace{u(2T)}_{u_2}\delta(t-2T)
  +\cdots+\underbrace{u(kT)}_{u_k}\delta(t-kT)+\cdots
\end{equation}
where $\sum_{k=0}^{\infty}\delta(t-kT)$ is the train of unit impulses.

\paragraph{Note 1: two alternative definitions of the $z$ transform.}
\begin{enumerate}
  \item[(1)] \emph{(old)} Given the sequence $\{u_0,u_1,\dots,u_k,\dots\}$, its $z$ transform is
  \begin{equation}
    U(z)=u_0+u_1z^{-1}+u_2z^{-2}+\cdots+u_kz^{-k}+\cdots=\sum_{k=0}^{\infty}u_kz^{-k}\qquad(\text{for }T=1),
  \end{equation}
  where $u_0=u(0)$, $u_1=u(T)$, \dots
  \item[(2)] Given the impulse-train representation of a discrete-time signal,
  $u^{*}(t)=\sum_{k=0}^{\infty}u_k\,\delta(t-kT)$, its Laplace transform is
  \begin{equation}
    U^{*}(s)=u_0+u_1\ee^{-sT}+u_2\ee^{-2sT}+\cdots+u_k\ee^{-ksT}+\cdots=\sum_{k=0}^{\infty}u_k\,\ee^{-ksT}.
  \end{equation}
  This is the \emph{alternative} $z$-transform definition: let $z$ be defined by
  \begin{equation}\label{eq:zes}
    z=\ee^{sT},
  \end{equation}
  so that $z$ and $s$ are related by this equation and $U^{*}(s)\big|_{z=\ee^{sT}}=U(z)$.
\end{enumerate}

\paragraph{Note 2: selection of the sampling period $T$} is very important.
\begin{theorem}[Sampling theorem]
If $\ws$ (the sampling frequency), defined as $\ws=\dfrac{2\pi}{T}$, satisfies $\ws>2\omega_1$, where
$\omega_1$ is the highest frequency component in the continuous-time signal $u(t)$, then $u(t)$ can be
reconstructed completely from the sampled signal $u^{*}(t)$.
\end{theorem}
In practice a rule of thumb is to choose $\ws=k\omega_1$ with $5\le k\le10$; the choice of the
constant $k$ depends on the application.

\begin{boardnote}
The board marks the range $5\le k\le10$ with a question mark. It is kept as written. Read it as a
starting point, not a law: how large $k$ must be depends on the application, and it is usually referred
to the fastest dynamics that matter (for control, the closed-loop bandwidth) rather than only to the
highest frequency in the signal.
\end{boardnote}

\addedsection[Added problems: sampling]
\begin{problem}[Choosing $T$]
A signal contains components up to $50\,\mathrm{Hz}$. (a) What is the largest sampling period allowed by
the sampling theorem? (b) What period results from the rule of thumb $5\le k\le10$?
\end{problem}
\begin{solution}
$\omega_1=2\pi\cdot50=314.16\,\mathrm{rad/s}$. (a) $\ws>2\omega_1=628.3\,\mathrm{rad/s}$, so
$T<2\pi/628.3=10\,\mathrm{ms}$. (b) $T=2\pi/(k\omega_1)=1/(50k)$ seconds, i.e.\ $4\,\mathrm{ms}$ for
$k=5$ and $2\,\mathrm{ms}$ for $k=10$.
\end{solution}

\begin{problem}[Aliasing]
A $80\,\mathrm{Hz}$ sine is sampled at $100\,\mathrm{Hz}$. At what frequency does the sampled sequence
appear to oscillate?
\end{problem}
\begin{solution}
Here $\ws=2\pi\cdot100$ is smaller than $2\cdot2\pi\cdot80$, so the sampling theorem is violated. The
sample values equal those of a sine at $|80-100|=20\,\mathrm{Hz}$ (with a sign change of phase), so
the $80\,\mathrm{Hz}$ tone is seen as a $20\,\mathrm{Hz}$ one.
\end{solution}

%------------------------------------------------------------------------------
\subsection{DAC model: the zero-order hold}\scanref{p.\,29}
\paragraph{Assumptions.}
\begin{itemize}
  \item The DAC outputs are exactly equal in magnitude to its inputs.
  \item The DAC yields an analog output instantaneously.
  \item The DAC outputs are constant over each sampling period.
\end{itemize}
Under these assumptions the input--output relationship of the DAC, for the input sequence $\{u(k)\}$, is
\begin{equation}
  u(t)=u(k),\qquad kT\le t<(k+1)T,\qquad k=0,1,2,\dots
\end{equation}
This is the \emph{zero-order hold} (ZOH).

\begin{figure}[ht]
\centering
\begin{tikzpicture}[x=1cm,y=1cm]
  % input sample
  \draw[thick] (-0.2,0) -- (1.2,0);
  \draw[thick,itublue] (0.5,0) -- (0.5,1.0); \fill[itublue] (0.5,1.0) circle (1.6pt);
  \node[lab,below] at (0.5,0) {$kT$};
  \draw[sig] (1.3,0.3) -- (2.0,0.3);
  \node[blk,minimum width=26mm,minimum height=14mm] (z) at (3.4,0.3) {Zero-order\\ hold};
  \draw[sig] (4.8,0.3) -- (5.5,0.3);
  % output pulse
  \draw[thick] (5.6,0) -- (8.6,0) node[right,lab]{$t$};
  \draw[thick,itublue] (5.9,0) -- (5.9,1.0) -- (7.4,1.0) -- (7.4,0);
  \draw[dotted] (5.9,0) -- (5.9,-0.15); \draw[dotted] (7.4,0) -- (7.4,-0.15);
  \node[lab,below] at (5.9,-0.1) {$kT$}; \node[lab,below] at (7.4,-0.1) {$(k{+}1)T$};
\end{tikzpicture}\\[6mm]
\begin{tikzpicture}[x=1cm,y=1cm]
  \draw[thick,->] (0,0) -- (7.4,0) node[right,lab]{$t$};
  \draw[thick,itublue] (0.6,0.8) -- (0.6,0) -- (0.6,0.8) -- (6.9,0.8);
  \draw[thick,itublue] (3.1,-0.8) -- (6.9,-0.8);
  \draw[dashed,gray] (0.6,0.8) -- (0.6,-0.1); \draw[dashed,gray] (3.1,0.8) -- (3.1,-0.8);
  \draw[thick,itublue] (0.6,0) -- (0.6,0.8);
  \draw[thick,itublue] (3.1,0) -- (3.1,-0.8);
  \node[lab,below left] at (0.6,0) {$kT$}; \node[lab,above right] at (3.1,0) {$(k{+}1)T$};
  \node[lab,above left,align=right] at (6.9,0.85) {positive step};
  \node[lab,below left,align=right] at (6.9,-0.85) {negative step};
\end{tikzpicture}
\caption{Zero-order hold: the held pulse equals a positive unit step at $kT$ plus a negative unit step
at $(k+1)T$ (scaled by $u(k)$).}
\end{figure}

\paragraph{ZOH transfer function.} Use the Laplace transform of the unit positive and negative steps and
the time-delay theorem:
\[
  \LT{1(t)}=\frac1s,\qquad \LT{1(t-T)}=\frac{\ee^{-sT}}{s}.
\]
Subtracting them (the impulse response of the ZOH is the unit pulse $1(t)-1(t-T)$) gives
\begin{equation}\label{eq:gzoh}
  \boxed{\,\GZOH(s)=\frac{1-\ee^{-sT}}{s}\,}
\end{equation}

\addedsection[Added problem: frequency response of the ZOH]
\begin{problem}
Show that $\GZOH(\jj\omega)=T\,\ee^{-\jj\omega T/2}\,\dfrac{\sin(\omega T/2)}{\omega T/2}$ and evaluate its
magnitude and phase at $\omega=\ws/2$.
\end{problem}
\begin{solution}
$\GZOH(\jj\omega)=\dfrac{1-\ee^{-\jj\omega T}}{\jj\omega}
=\ee^{-\jj\omega T/2}\dfrac{\ee^{\jj\omega T/2}-\ee^{-\jj\omega T/2}}{\jj\omega}
=T\,\ee^{-\jj\omega T/2}\dfrac{\sin(\omega T/2)}{\omega T/2}$. At $\omega=\ws/2=\pi/T$ we get
$\omega T/2=\pi/2$: the magnitude is $T\cdot\dfrac{1}{\pi/2}=0.637\,T$ and the phase is $-\pi/2=-90^{\circ}$. The ZOH therefore
behaves like a low-pass filter with a delay of $T/2$ seconds.
\end{solution}

%------------------------------------------------------------------------------
\subsection{The transfer function $\GZAS(z)$}\scanref{p.\,30}
The transfer function of the DAC and the analog subsystem in cascade is
\begin{equation}
  G_{\mathrm{ZA}}(s)=G(s)\,\GZOH(s)=\left(\frac{1-\ee^{-sT}}{s}\right)G(s)=(1-\ee^{-sT})\frac{G(s)}{s}.
\end{equation}
The corresponding impulse response is
\begin{equation}
  g_{\mathrm{ZA}}(t)=g(t)*g_{\mathrm{ZOH}}(t)=g_s(t)-g_s(t-T),
\end{equation}
\[
  g_s(t)=\mathcal{L}^{-1}\!\left[\frac{G(s)}{s}\right],\qquad
  g_s(t-T)=\mathcal{L}^{-1}\!\left[\ee^{-sT}\frac{G(s)}{s}\right].
\]
\begin{itemize}
  \item The impulse response is the analog system's \emph{step response} minus a second step response
        delayed by one sampling period.
  \item It is sampled to give the sampled impulse response
        \[
          g_{\mathrm{ZA}}(kT)=g_s(kT)-g_s(kT-T).
        \]
\end{itemize}
By $z$ transforming we obtain the $z$ transfer function of the DAC (ZOH), the analog subsystem and the
ideal sampler:
\begin{align}
  \GZAS(z)&=\ZT{g_s^{*}(t)}-\ZT{g_s^{*}(t-T)}
  =\ZT{\left[\mathcal{L}^{-1}\!\left[\tfrac{G(s)}{s}\right]\right]^{*}}
   -z^{-1}\,\ZT{\left[\mathcal{L}^{-1}\!\left[\tfrac{G(s)}{s}\right]\right]^{*}}\notag\\
  &=(1-z^{-1})\,\ZT{\left[\mathcal{L}^{-1}\!\left[\tfrac{G(s)}{s}\right]\right]^{*}}.
\end{align}
The star is written to emphasize that sampling of a time function is necessary before the $z$ transform.
Finally,
\begin{equation}\label{eq:gzas}
  \boxed{\,\GZAS(z)=(1-z^{-1})\,\ZT{\frac{G(s)}{s}}\,}
\end{equation}
In calculations we use tables that go directly from $s$ to $z$ ($s\to t\to k\to z$).

\begin{remark}[added]
$\GZAS(z)$ reproduces, at the sampling instants, the response of the analog system to a piecewise-constant
input. In particular its DC gain is that of the analog system: $\GZAS(1)=G(0)$ (see the examples below).
\end{remark}

%------------------------------------------------------------------------------
\subsection{Examples of $\GZAS(z)$}\scanref{pp.\,31--34}

\begin{example}[Cruise (velocity) control of a vehicle]
Find $\GZAS(z)$ for the cruise control system of a vehicle, where $u$ is the input force, $v$ is the
velocity of the car and $b$ is the viscous friction coefficient (Fig.~\ref{fig:cruise}).

\begin{figure}[H]
\centering
\begin{tikzpicture}[x=1cm,y=1cm]
  % car
  \node[draw=itublue,thick,minimum width=16mm,minimum height=9mm,fill=white] (m) at (0,0) {$M$};
  \fill[itublue!20,draw=itublue,thick] (-0.5,-0.55) circle (0.17); \fill[itublue!20,draw=itublue,thick] (0.5,-0.55) circle (0.17);
  \draw[thick] (-1.6,-0.72) -- (1.6,-0.72);
  \draw[sig] (m.east) -- ++(1.1,0) node[right,lab]{$u(t)$};
  \draw[sig,boardred] (m.west) -- ++(-1.1,0) node[left,lab,text=boardred]{$b\,v(t)$};
  \draw[sig] (-0.5,0.75) -- (0.8,0.75) node[right,lab]{$x,\,v$};
\end{tikzpicture}
\caption{Vehicle model for the cruise-control example.}\label{fig:cruise}
\end{figure}

\noindent The control loop is the one of Fig.~\ref{fig:digitalloop} with $R(z)=V_{\mathrm{ref}}(z)$ and output $V(z)$.
Using Newton's law (initial conditions zero):
\[
  M\dot v(t)=u(t)-b\,v(t)\ \Rightarrow\ Ms\,V(s)=U(s)-b\,V(s),\qquad
  G(s)=\frac{V(s)}{U(s)}=\frac{1}{Ms+b}.
\]
Rewrite the function as
\[
  G(s)=\frac{K}{\tau s+1}=\frac{K/\tau}{s+1/\tau}=\frac{1/M}{s+b/M},\qquad \tau=\frac Mb,\quad K=\frac1b.
\]
Then $\GZAS(z)=(1-z^{-1})\,\ZT{G(s)/s}$ with the partial-fraction expansion
\[
  \frac{G(s)}{s}=\frac1s\left(\frac{K/\tau}{s+1/\tau}\right)=\frac{A_1}{s}+\frac{A_2}{s+1/\tau},
\]
\[
  A_1=\left[s\cdot\frac1s\frac{K/\tau}{s+1/\tau}\right]_{s=0}=K,\qquad
  A_2=\left[\left(s+\tfrac1\tau\right)\frac1s\frac{K/\tau}{s+1/\tau}\right]_{s=-1/\tau}=-K,
\]
\[
  \Rightarrow\quad \frac{G(s)}{s}=\frac Ks-\frac{K}{s+1/\tau}.
\]
Using $\ZT{1/s}=\dfrac{1}{1-z^{-1}}$ and $\ZT{1/(s+a)}=\dfrac{1}{1-\ee^{-aT}z^{-1}}$:
\begin{align*}
  \GZAS(z)&=(1-z^{-1})\left[\frac{K}{1-z^{-1}}-\frac{K}{1-\ee^{-T/\tau}z^{-1}}\right]
  =K-\frac{K(1-z^{-1})}{1-\ee^{-T/\tau}z^{-1}}\\
  &=\frac{K\bigl(1-\ee^{-T/\tau}z^{-1}-1+z^{-1}\bigr)}{1-\ee^{-T/\tau}z^{-1}}
  =\frac{K z^{-1}\bigl(1-\ee^{-T/\tau}\bigr)}{1-\ee^{-T/\tau}z^{-1}}\\
  &=\boxed{\ \frac{K\bigl(1-\ee^{-T/\tau}\bigr)}{z-\ee^{-T/\tau}}\ }
\end{align*}
\end{example}

\begin{example}[Position control of a vehicle]
Find $\GZAS(z)$ for the position of the car, with $u$ the input force, $x$ the position and $b$ the viscous
friction coefficient. Newton's law gives
\[
  M\ddot x(t)=u(t)-b\,\dot x(t)\ \Rightarrow\ Ms^{2}X(s)=U(s)-bs\,X(s),\qquad
  G(s)=\frac{X(s)}{U(s)}=\frac{1}{s(Ms+b)}.
\]
Rewrite the function ($\tau=M/b$, $K=1/b$):
\[
  G(s)=\frac{K}{s(\tau s+1)}=\frac{K/\tau}{s(s+1/\tau)},\qquad
  \frac{G(s)}{s}=\frac{K/\tau}{s^{2}(s+1/\tau)}=\frac{A_{11}}{s^{2}}+\frac{A_{12}}{s}+\frac{A_2}{s+1/\tau}.
\]
The coefficients are
\[
  A_{11}=\left[s^{2}\,\frac{K/\tau}{s^{2}(s+1/\tau)}\right]_{s=0}=K,
\]
\[
  A_{12}=\frac{1}{1!}\frac{\mathrm d}{\mathrm ds}\left[\frac{K/\tau}{s+1/\tau}\right]_{s=0}
  =\frac K\tau\left[\frac{-1}{(s+1/\tau)^{2}}\right]_{s=0}=-K\tau,
\]
\[
  A_2=\left[\left(s+\tfrac1\tau\right)\frac{K/\tau}{s^{2}(s+1/\tau)}\right]_{s=-1/\tau}=K\tau,
\]
\[
  \Rightarrow\qquad \frac{G(s)}{s}=\frac{K}{s^{2}}-\frac{K\tau}{s}+\frac{K\tau}{s+1/\tau}.
\]
\end{example}

\begin{boardnote}
On the board the $z$ transform of $K/s^{2}$ is written as $Kz/(z-1)^{2}$, i.e.\ with the factor $T$ left
out (valid for $T=1\,\mathrm{s}$). The general result uses $\ZT{1/s^{2}}=Tz/(z-1)^{2}$ and is derived below;
for $T=1$ it reduces exactly to the final line of the board.
\end{boardnote}

Continuing the position example for general $T$, with $\alpha=\ee^{-T/\tau}$ and
$\ZT{1/s^{2}}=\dfrac{Tz}{(z-1)^{2}}$:
\begin{align*}
  \GZAS(z)&=\frac{z-1}{z}\left[\frac{KTz}{(z-1)^{2}}-\frac{K\tau z}{z-1}+\frac{K\tau z}{z-\alpha}\right]
  =\frac{KT}{z-1}-K\tau+\frac{K\tau(z-1)}{z-\alpha}\\
  &=K\,\frac{T(z-\alpha)-\tau(z-1)(z-\alpha)+\tau(z-1)^{2}}{(z-1)(z-\alpha)}\\
  &=K\,\frac{(T-\tau+\tau\alpha)\,z+(\tau-\tau\alpha-T\alpha)}{(z-1)(z-\alpha)}.
\end{align*}
\begin{equation}\label{eq:position-zas}
  \boxed{\ \GZAS(z)=K\,\frac{(T-\tau+\tau\alpha)\,z+\bigl(\tau-\alpha(\tau+T)\bigr)}{(z-1)(z-\alpha)},\qquad\alpha=\ee^{-T/\tau}\ }
\end{equation}
For $T=1\,\mathrm{s}$ this is the board result
$K\,\dfrac{(1-\tau+\tau\ee^{-1/\tau})z+(\tau-\ee^{-1/\tau}(\tau+1))}{(z-1)(z-\ee^{-1/\tau})}$.

\begin{example}[Series R--L circuit, inductor voltage as output]\label{ex:rl}
Find $\GZAS(z)$ for the series R--L circuit with the inductor voltage as output.

\begin{figure}[H]
\centering
\begin{circuitikz}[european resistors,european inductors,scale=0.9,transform shape]
  \draw (0,2) to[R=$R$,o-] (3,2) to[short,-o] (4.6,2);
  \draw (3,2) to[L=$L$] (3,0);
  \draw (0,0) to[short,o-o] (4.6,0);
  \fill (3,2) circle (1.6pt); \fill (3,0) circle (1.6pt);
  \draw[<->,thick] (-0.35,0.15) -- (-0.35,1.85) node[midway,left]{$V_{in}$};
  \draw[<->,thick] (4.95,0.15) -- (4.95,1.85) node[midway,right]{$V_o$};
\end{circuitikz}
\caption{Series R--L circuit.}
\end{figure}

\noindent With the current $i$:
\[
  V_{in}=Ri+L\frac{\mathrm di}{\mathrm dt},\qquad V_o=L\frac{\mathrm di}{\mathrm dt}
\]
\[
  G(s)=\frac{V_o(s)}{V_{in}(s)}=\frac{Ls\,I(s)}{Ls\,I(s)+R\,I(s)}=\frac{Ls}{Ls+R}=\frac{s}{s+1/\tau},
\]
where $\tau=L/R$.
\[
  \GZAS(z)=(1-z^{-1})\,\ZT{\frac{G(s)}{s}}=(1-z^{-1})\,\ZT{\frac{1}{s+1/\tau}}
  =\frac{z-1}{z}\cdot\frac{z}{z-\ee^{-T/\tau}}=\boxed{\ \frac{z-1}{z-\ee^{-T/\tau}}\ }
\]
\end{example}

\addedsection[Added problems: computing $\GZAS(z)$]
\begin{problem}[Cruise control with numbers]
A car has mass $M=1000\,\mathrm{kg}$ and friction coefficient $b=50\,\mathrm{N\,s/m}$. Find $\GZAS(z)$ for
$T=1\,\mathrm{s}$, and check its DC gain.
\end{problem}
\begin{solution}
$\tau=M/b=20\,\mathrm{s}$, $K=1/b=0.02$, $\ee^{-T/\tau}=\ee^{-0.05}=0.951229$. Hence
\[
  \GZAS(z)=\frac{0.02\,(1-0.951229)}{z-0.951229}=\frac{0.000975}{z-0.9512}.
\]
DC gain: $\GZAS(1)=\dfrac{0.000975}{1-0.9512}=0.02=K=G(0)$. \;(Checked against a ZOH discretization with
\texttt{scipy.signal.cont2discrete}.)
\end{solution}

\begin{problem}[Position control with numbers]
Evaluate Eq.~\eqref{eq:position-zas} for $M=b=1$ (so $\tau=K=1$) and $T=0.5\,\mathrm{s}$.
\end{problem}
\begin{solution}
$\ee^{-0.5}=0.606531$. Numerator: $(T-\tau+\tau\ee^{-0.5})=0.5-1+0.6065=0.1065$ and
$\tau-\ee^{-0.5}(\tau+T)=1-0.6065\cdot1.5=0.0902$. Therefore
\[
  \GZAS(z)=\frac{0.1065\,z+0.0902}{(z-1)(z-0.6065)}=\frac{0.1065\,z+0.0902}{z^{2}-1.6065\,z+0.6065}.
\]
Both coefficients agree with a numerical ZOH discretization of $1/(s(s+1))$.
\end{solution}

%------------------------------------------------------------------------------
\subsection{Starred Laplace transform and the pulse transfer function}\scanref{pp.\,34--35}
In analyzing discrete-time control systems we often find that some signals are \emph{starred}, meaning that
the signals are impulse sampled:
\begin{itemize}
  \item $x(t)$: continuous-time signal,
  \item $x^{*}(t)$: impulse-sampled signal (with an ideal sampler).
\end{itemize}

\begin{figure}[ht]
\centering
\resizebox{0.8\linewidth}{!}{\figStarred}
\caption{Impulse-sampled input and output of $G(s)$.}
\end{figure}

\paragraph{Pulse transfer function.} From the figure, $Y(s)=G(s)X^{*}(s)$. Taking the starred transform
(we may factor out $X^{*}(s)$),
\begin{equation}
  Y^{*}(s)=\bigl[G(s)X^{*}(s)\bigr]^{*}=G^{*}(s)\,X^{*}(s).
\end{equation}
\begin{equation}
  G^{*}(s)=\frac{Y^{*}(s)}{X^{*}(s)}\qquad\text{(the pulse transfer function).}
\end{equation}
The signal chain is: inverse Laplace transform $\to y(t)$; impulse sampling $\to y^{*}(t),\ y(kT)$;
Laplace transform and $z$ transform $\to Y^{*}(s),\ Y(z)$.

\begin{itemize}
  \item $G(s)$: ordinary transfer function, non-periodic.
  \item $G^{*}(s)$: pulse transfer function, \emph{periodic} (it depends on $\ee^{sT}$, $T=2\pi/\ws$).
  \item Such periodic systems are generally more complicated to analyze than the original non-periodic
        system, but they may be analyzed without difficulty if the analysis is carried out in the
        $z$ domain: $G^{*}(s)\leftrightarrow G(z)$ with $z=\ee^{sT}$.
\end{itemize}

\addedsection[Added problem: why $X^{*}(s)$ factors out]
\begin{problem}
Show that $Y^{*}(s)=G^{*}(s)X^{*}(s)$ when $Y(s)=G(s)X^{*}(s)$.
\end{problem}
\begin{solution}
In the time domain $y(t)=\sum_{m=0}^{\infty}x(mT)\,g(t-mT)$, so at the sampling instants
$y(kT)=\sum_{m=0}^{k}x(mT)\,g\bigl((k-m)T\bigr)$ (a discrete convolution, since $g(t)=0$ for $t<0$).
The $z$ transform of a discrete convolution is the product of the $z$ transforms:
$Y(z)=G(z)X(z)$, and with $z=\ee^{sT}$ this is $Y^{*}(s)=G^{*}(s)X^{*}(s)$.
\end{solution}

%------------------------------------------------------------------------------
\subsection{Effect of the sampler on the transfer function of a cascade}\scanref{pp.\,36--38}
In a discrete-time system that includes several analog subsystems in cascade and several samplers, the
\emph{location of the samplers} plays an important role in determining the overall transfer function.

\paragraph{\textcircled{1} The cascade is separated by samplers.} Each block then has a sampled output and
input, as well as a $z$-domain transfer function.
\begin{center}
\resizebox{0.9\linewidth}{!}{\figCascadeS}
\end{center}
\begin{align*}
  X(s)&=H_1(s)\,U^{*}(s), & X^{*}(s)&=H_1^{*}(s)\,U^{*}(s),\\
  Y(s)&=H_1^{*}(s)\,H_2(s)\,U^{*}(s), & Y^{*}(s)&=H_1^{*}(s)\,H_2^{*}(s)\,U^{*}(s)
  \ \xrightarrow{\ z=\ee^{sT}\ }\ Y(z)=H_1(z)\,H_2(z)\,U(z).
\end{align*}

\paragraph{\textcircled{2} The cascade is not separated by samplers.} The components cannot be separated
after sampling.
\begin{center}
\resizebox{0.78\linewidth}{!}{\figCascadeNS}
\end{center}
\begin{align*}
  Y(s)&=\bigl[H_1(s)H_2(s)\bigr]U^{*}(s),\\
  Y^{*}(s)&=\bigl[H_1(s)H_2(s)\bigr]^{*}U^{*}(s)=(H_1H_2)^{*}(s)\,U^{*}(s)
  \ \xrightarrow{\ z=\ee^{sT}\ }\ Y(z)=(H_1H_2)(z)\,U(z).
\end{align*}
\textbf{Note:} the overall transfer functions are different in the two cases:
$H_1(z)H_2(z)\ne(H_1H_2)(z)$, where $(H_1H_2)(z)=\ZT{H_1(s)H_2(s)}$.

\begin{example}[Cascade of two analog systems with sampled input]
Find the equivalent sampled impulse response sequence $h(kT)$ and the equivalent $z$-transfer function of the
cascade of the two analog systems $H_1(s)=\dfrac{1}{s+2}$ and $H_2(s)=\dfrac{2}{s+4}$ with sampled input.

\begin{center}
\resizebox{0.78\linewidth}{!}{\figCascadeNS}\\[1mm]
{\small 1) direct (no sampler between the blocks)}\\[3mm]
\resizebox{0.9\linewidth}{!}{\figCascadeS}\\[1mm]
{\small 2) separated (a sampler between the blocks)}
\end{center}

\textbf{1) No sampler between the blocks.}
\[
  H(s)=H_1(s)H_2(s)=\frac{2}{(s+2)(s+4)}=\frac{1}{s+2}-\frac{1}{s+4}.
\]
The impulse response of this cascade system is $h(t)=\ee^{-2t}-\ee^{-4t}$; with $t=kT$,
\[
  h(kT)=\ee^{-2kT}-\ee^{-4kT},\qquad k\ge0,
\]
and thus the $z$ transform is
\[
  H(z)=\frac{z}{z-\ee^{-2T}}-\frac{z}{z-\ee^{-4T}}=\frac{(\ee^{-2T}-\ee^{-4T})\,z}{(z-\ee^{-2T})(z-\ee^{-4T})}.
\]

\textbf{2) A sampler between the blocks.}
\[
  H_1(s)=\frac{1}{s+2}\xrightarrow{\ z\ }H_1(z)=\frac{z}{z-\ee^{-2T}},\qquad
  H_2(s)=\frac{2}{s+4}\xrightarrow{\ z\ }H_2(z)=\frac{2z}{z-\ee^{-4T}},
\]
\[
  H(z)=H_1(z)H_2(z)=\frac{2z^{2}}{(z-\ee^{-2T})(z-\ee^{-4T})}
\]
\[
  \Rightarrow\ H(z)=\frac{2}{\ee^{-2T}-\ee^{-4T}}\left(\frac{\ee^{-2T}z}{z-\ee^{-2T}}-\frac{\ee^{-4T}z}{z-\ee^{-4T}}\right).
\]
Inverse $z$ transform:
\[
  h(kT)=\frac{2}{\ee^{-2T}-\ee^{-4T}}\Bigl[\ee^{-2T}\ee^{-2kT}-\ee^{-4T}\ee^{-4kT}\Bigr],\qquad k\ge0.
\]
\end{example}

The two answers have the same poles ($\ee^{-2T}$ and $\ee^{-4T}$) but different numerators: the first has
a single zero at $z=0$ and the second a double zero at $z=0$.

\addedsection[Added problem: the sampler position matters]
\begin{problem}
For $H_1(s)=\dfrac{1}{s+1}$, $H_2(s)=\dfrac{1}{s+2}$ and $T=0.5\,\mathrm{s}$ find the two overall transfer
functions (no sampler between the blocks, and a sampler between the blocks), and compare their first four
samples of the impulse response.
\end{problem}
\begin{solution}
With $a=\ee^{-0.5}=0.6065$ and $b=\ee^{-1}=0.3679$ the denominator is common:
$(z-a)(z-b)=z^{2}-0.9744\,z+0.2231$.\par
\emph{No sampler between the blocks:} $H_1H_2=\dfrac{1}{s+1}-\dfrac{1}{s+2}$, so
$H(z)=\dfrac{(a-b)\,z}{(z-a)(z-b)}=\dfrac{0.2387\,z}{z^{2}-0.9744z+0.2231}$ and
$h(kT)=\ee^{-0.5k}-\ee^{-k}=0,\;0.2387,\;0.2325,\;0.1733$.\par
\emph{Sampler between the blocks:} $H(z)=\dfrac{z}{z-a}\cdot\dfrac{z}{z-b}=\dfrac{z^{2}}{z^{2}-0.9744z+0.2231}$ and
$h(kT)=1,\;0.9744,\;0.7263,\;0.4903$ (the discrete convolution of $\ee^{-0.5k}$ and $\ee^{-k}$).
The two transfer functions differ even though the analog blocks are identical.
\end{solution}

%------------------------------------------------------------------------------
\subsection{The closed-loop transfer function}\scanref{p.\,38}
For the loop of Fig.~\ref{fig:closedloop} (controller $C(z)$, then DAC (ZOH), analog subsystem and ADC
(ideal sampler), which together are $\GZAS(z)$):

\begin{figure}[ht]
\centering
\resizebox{\linewidth}{!}{\figClosedLoopExpanded}
\caption{Closed-loop digital control system; the box $\GZAS(z)$ contains the ZOH, the analog
subsystem and the ideal sampler.}\label{fig:closedloop}
\end{figure}

\begin{equation}
  G_{cl}(z)=\frac{Y(z)}{R(z)}=\frac{C(z)\,\GZAS(z)}{1+C(z)\,\GZAS(z)},
  \qquad\text{characteristic equation: }\ 1+C(z)\,\GZAS(z)=0.
\end{equation}
The roots of the characteristic equation are the closed-loop system poles, which can be selected for the
desired time-response specifications as in $s$-domain design.

\addedsection[Added problem: closed loop with the cruise-control plant]
\begin{problem}
Use $\GZAS(z)=\dfrac{0.000975}{z-0.9512}$ from the cruise-control problem ($T=1\,\mathrm{s}$) with
$C(z)=K_c$. (a) Find the closed-loop pole as a function of $K_c$ and the range of $K_c$ that gives a
stable loop. (b) Choose $K_c$ to put the pole at $z=0.5$ and compute $G_{cl}(1)$.
\end{problem}
\begin{solution}
(a) $1+K_c\GZAS(z)=0\Rightarrow z=0.9512-0.000975\,K_c$. Stability requires $|z|<1$, i.e.\
$-1<0.9512-0.000975K_c<1$, which gives $0<K_c<(1+0.9512)/0.000975\approx2000$.\par
(b) $0.9512-0.000975K_c=0.5\Rightarrow K_c=462.6$. With $K=0.02$,
$G_{cl}(1)=\dfrac{K_cK}{1+K_cK}=\dfrac{9.252}{10.252}=0.902$: the closed-loop DC gain is below one (a
proportional controller leaves a steady-state error of $9.8\%$ for a step command; this is the topic of the
next chapter).
\end{solution}

%------------------------------------------------------------------------------
\subsection{Example: sampled feedback with an analog loop}\scanref{p.\,39}
\begin{example}
Find the Laplace transform of the analog output $Y(s)$ and of the sampled output $Y^{*}(s)$ for the
block diagram of Fig.~\ref{fig:sampledfb}.

\begin{figure}[H]
\centering
\resizebox{0.92\linewidth}{!}{\figSampledFeedback}
\caption{Analog loop with the feedback signal sampled.}\label{fig:sampledfb}
\end{figure}

\noindent From the diagram,
\[
  X(s)=H(s)Y(s),\qquad E(s)=R(s)-X^{*}(s),\qquad Y(s)=G(s)D(s)E(s),
\]
so that
\[
  X(s)=H(s)G(s)D(s)\bigl[R(s)-X^{*}(s)\bigr]=H(s)G(s)D(s)R(s)-H(s)G(s)D(s)X^{*}(s).
\]
The impulse-sampled signal $x^{*}(t)$ has the Laplace form
\[
  X^{*}(s)=(HGDR)^{*}(s)-(HGD)^{*}(s)X^{*}(s)\quad\longrightarrow\quad X^{*}(s)=\frac{(HGDR)^{*}(s)}{1+(HGD)^{*}(s)}.
\]
Hence
\[
  E(s)=R(s)-\frac{(HGDR)^{*}(s)}{1+(HGD)^{*}(s)},\qquad
  Y(s)=G(s)D(s)\left[R(s)-\frac{(HGDR)^{*}(s)}{1+(HGD)^{*}(s)}\right].
\]
The sampled output is
\[
  Y^{*}(s)=(GDR)^{*}(s)-(GD)^{*}(s)\,\frac{(HGDR)^{*}(s)}{1+(HGD)^{*}(s)}.
\]
With the transformation $z=\ee^{sT}$ the sampled output in the $z$ domain is
\[
  Y(z)=(GDR)(z)-(GD)(z)\,\frac{(HGDR)(z)}{1+(HGD)(z)}.
\]
\end{example}
\begin{itemize}
  \item This shows that for some digital systems no expression is available for the transfer function
        excluding the input $R$ (the input stays inside the $z$ transforms and cannot be factored out).
  \item However, we do have the characteristic equation $1+(HGD)(z)=0$ with which to design a control
        system, as in the cases where a closed-loop transfer function is defined.
\end{itemize}

%------------------------------------------------------------------------------
\subsection{Analog disturbances in a digital system}\scanref{pp.\,40--41}
\begin{itemize}
  \item Disturbances are variables that are not included in the system model but affect its response.
  \item Almost all disturbances are analog and are inputs to analog subsystems.
\end{itemize}

\begin{figure}[ht]
\centering
\resizebox{\linewidth}{!}{\figDisturbance}
\caption{Digital control loop with an analog disturbance $D(s)$ entering through the disturbance model
$G_d(s)$.}\label{fig:disturbance}
\end{figure}

From Fig.~\ref{fig:disturbance},
\[
  Y(s)=G(s)\bigl[G_d(s)D(s)+\GZOH(s)U^{*}(s)\bigr],\qquad U(s)=C^{*}(s)E(s),\qquad E(s)=R(s)-Y^{*}(s),
\]
so $U(s)=C^{*}(s)\bigl[R(s)-Y^{*}(s)\bigr]$. The system is linear, so the reference input can be treated
separately and is assumed to be zero: $R(s)=0$. Then
\[
  U(s)=C^{*}(s)\bigl[0-Y^{*}(s)\bigr]=-C^{*}(s)Y^{*}(s)\ \longrightarrow\ U^{*}(s)=-C^{*}(s)Y^{*}(s),
\]
\[
  Y(s)=G(s)\bigl[G_d(s)D(s)-\GZOH(s)C^{*}(s)Y^{*}(s)\bigr],
\]
\[
  Y^{*}(s)=(GG_dD)^{*}(s)-(G\,\GZOH)^{*}(s)\,C^{*}(s)Y^{*}(s)
  \ \Rightarrow\ Y^{*}(s)=\frac{(GG_dD)^{*}(s)}{1+\GZAS^{*}(s)\,C^{*}(s)},
\]
where $G\,\GZOH$ is the $\GZAS$ block. In the $z$ domain,
\begin{equation}\label{eq:dist}
  Y(z)=\frac{(GG_dD)(z)}{1+\GZAS(z)\,C(z)}.
\end{equation}

\begin{example}[Steady-state response to an impulse disturbance]
Consider the block diagram of a digital system with an analog disturbance and
\[
  G(s)=\frac{K_p}{s+1},\qquad G_d(s)=\frac1s,\qquad C(z)=K_c,\qquad y(\infty)=\ ?
\]
Find the steady-state response of the system to an impulse disturbance of magnitude $A$. From
Eq.~\eqref{eq:dist}, $Y(z)=\dfrac{(GG_dD)(z)}{1+\GZAS(z)C(z)}$ with $\GZAS(z)=(1-z^{-1})\ZT{G(s)/s}$.
For the disturbance, $D(s)=A$ and
\[
  G(s)G_d(s)D(s)=\frac{K_pA}{s(s+1)}=K_pA\left(\frac1s-\frac{1}{s+1}\right),
\]
\[
  (GG_dD)(z)=K_pA\left[\frac{z}{z-1}-\frac{z}{z-\ee^{-T}}\right].
\]
\[
  \GZAS(z)=(1-z^{-1})\,K_p\left(\frac{z}{z-1}-\frac{z}{z-\ee^{-T}}\right)=K_p\,\frac{1-\ee^{-T}}{z-\ee^{-T}}.
\]
\[
  Y(z)=\frac{K_pA\left[\dfrac{z}{z-1}-\dfrac{z}{z-\ee^{-T}}\right]}{1+K_cK_p\dfrac{1-\ee^{-T}}{z-\ee^{-T}}}.
\]
For all steady-state responses we use the final value theorem:
\[
  y(\infty)=\lim_{z\to1}\bigl[(1-z^{-1})Y(z)\bigr]=\boxed{\ \frac{K_pA}{1+K_cK_p}\ }
\]
\end{example}

\begin{remark}[added]
The final value theorem is valid only if the closed loop is stable. Here the closed-loop pole is
$z=\ee^{-T}-K_cK_p(1-\ee^{-T})$, which lies inside the unit circle for
$-1<\ee^{-T}-K_cK_p(1-\ee^{-T})<1$, i.e.\ for $K_cK_p<\dfrac{1+\ee^{-T}}{1-\ee^{-T}}$ (for $K_cK_p>0$).
\end{remark}

\addedsection[Added problem: numerical check of the disturbance example]
\begin{problem}
In the example above take $K_p=2$, $T=0.5\,\mathrm{s}$, $A=1$. Compute $y(\infty)$ for $K_c=0.25$, and
discuss $K_c=3.5$.
\end{problem}
\begin{solution}
For $K_c=0.25$: $K_cK_p=0.5$, $y(\infty)=\dfrac{2\cdot1}{1+0.5}=1.333$. The closed-loop pole is
$\ee^{-0.5}-0.5(1-\ee^{-0.5})=0.6065-0.1967=0.410$, inside the unit circle, so the formula applies (a
simulation of the loop converges to $1.3333$).\par
For $K_c=3.5$: $K_cK_p=7$ and the pole is $0.6065-7\cdot0.3935=-2.148$, outside the unit circle, so the
loop is unstable and $y(k)$ diverges although the formula would still produce the number
$\dfrac{2}{1+7}=0.25$. The limit $K_cK_p<\dfrac{1+\ee^{-0.5}}{1-\ee^{-0.5}}=4.08$ is violated.
\end{solution}
